\chapter{Introduction to Machine Learning}
\label{ch:introduction}

\epigraph{``Machine learning is the science of getting computers to act without being explicitly programmed.''}{Arthur Samuel, 1959}

\learninggoal{
	\item Define machine learning and distinguish it from traditional programming.
	\item Classify learning problems: supervised, unsupervised, reinforcement learning.
	\item Understand the bias-variance tradeoff as the central tension in ML.
	\item Formulate learning as function approximation from data.
}

\section{What Is Machine Learning?}

Traditional programming: a human writes explicit rules (a program) that maps inputs to outputs. Machine learning inverts this: the computer learns the rules from examples.

\begin{definition}[Machine Learning]
	A computer program is said to \textbf{learn} from experience $E$ with respect to some class of tasks $T$ and performance measure $P$, if its performance at tasks in $T$, as measured by $P$, improves with experience $E$.
\end{definition}

In practical terms: given a dataset $\mathcal{D} = \{(\mathbf{x}_i, y_i)\}_{i=1}^n$, find a function $f$ such that $f(\mathbf{x}) \approx y$ for new, unseen inputs.

\section{The Three Learning Paradigms}

\begin{longtable}{p{3cm}p{4.5cm}p{5cm}}
	\toprule
	\textbf{Paradigm}      & \textbf{Data}                         & \textbf{Goal}                 \\
	\midrule
	Supervised learning    & Input--output pairs $(\mathbf{x}, y)$ & Predict $y$ from $\mathbf{x}$ \\
	Unsupervised learning  & Inputs only $\{\mathbf{x}_i\}$        & Discover hidden structure     \\
	Reinforcement learning & States, actions, rewards              & Learn optimal policy          \\
	\bottomrule
\end{longtable}

This book focuses on supervised and unsupervised learning, which form the foundation of most practical ML.

\subsection{Supervised Learning}

Given labeled examples $\{(\mathbf{x}_i, y_i)\}_{i=1}^n$, learn a mapping $f: \X \to \Y$:

\begin{itemize}
	\item \textbf{Regression:} $\Y = \R$ (continuous output). Examples: house prices, temperature.
	\item \textbf{Classification:} $\Y = \{1, \dots, K\}$ (discrete classes). Examples: spam detection, image recognition.
\end{itemize}

The quality of $f$ is measured by a \emph{loss function} $L(y, f(\mathbf{x}))$:

\begin{itemize}
	\item Squared error for regression: $L(y, \hat{y}) = (y - \hat{y})^2$
	\item Cross-entropy for classification: $L(y, \hat{y}) = -\sum_{k} y_k \log \hat{y}_k$
\end{itemize}

\subsection{Unsupervised Learning}

Given unlabeled inputs $\{\mathbf{x}_i\}_{i=1}^n$, discover structure:

\begin{itemize}
	\item \textbf{Clustering:} group similar points (K-means, DBSCAN).
	\item \textbf{Dimensionality reduction:} find a low-dimensional representation (PCA).
	\item \textbf{Density estimation:} model $p(\mathbf{x})$ directly.
\end{itemize}

\section{The Central Tension: Bias--Variance Tradeoff}

Every model makes a tradeoff between two sources of error:

\begin{definition}[Bias and Variance]
	The \textbf{bias} of an estimator measures how far the average prediction is from the true value. The \textbf{variance} measures how much the prediction changes across different training sets.
\end{definition}

For a model $f$ trained on dataset $\mathcal{D}$:

\[
	\E_{\mathcal{D}}\big[(y - f(\mathbf{x}; \mathcal{D}))^2\big] = \underbrace{\big(\E_{\mathcal{D}}[f(\mathbf{x}; \mathcal{D})] - \E[y|\mathbf{x}]\big)^2}_{\text{Bias}^2} + \underbrace{\Var_{\mathcal{D}}\big(f(\mathbf{x}; \mathcal{D})\big)}_{\text{Variance}} + \underbrace{\sigma^2}_{\text{Irreducible error}}
\]

\begin{itemize}
	\item \textbf{High bias:} model is too simple, underfits the data (e.g., linear regression on a quadratic function).
	\item \textbf{High variance:} model is too complex, overfits the noise (e.g., high-degree polynomial).
\end{itemize}

The goal of feature engineering, regularization, and model selection is to find the sweet spot.

\begin{principle}[No Free Lunch]
	No single algorithm works best for all problems. Every model encodes assumptions (inductive bias). The art of ML is matching the bias to the data.
\end{principle}

\section{Learning as Function Approximation}

From a statistical perspective, all of supervised learning reduces to:

\[
	y = f^*(\mathbf{x}) + \epsilon, \quad \epsilon \sim \N(0, \sigma^2)
\]

where $f^*$ is the true (unknown) function and $\epsilon$ is irreducible noise. We choose a hypothesis class $\mathcal{H}$ (e.g., linear functions, decision trees) and search for the best $f \in \mathcal{H}$ that minimizes empirical risk:

\[
	\hat{f} = \arg\min_{f \in \mathcal{H}} \frac{1}{n} \sum_{i=1}^n L(y_i, f(\mathbf{x}_i))
\]

The gap between $\hat{f}$ and $f^*$ depends on:
\begin{enumerate}
	\item \textbf{Approximation error:} how well $\mathcal{H}$ can represent $f^*$.
	\item \textbf{Estimation error:} how well $\hat{f}$ approximates the best in $\mathcal{H}$ given finite data.
	\item \textbf{Optimization error:} how well our algorithm finds $\hat{f}$.
\end{enumerate}

\section{The Role of Feature Engineering}

No amount of statistical sophistication compensates for poor features. Feature engineering---transforming raw data into informative predictors---is often the highest-leverage activity in applied ML.

\begin{quote}
	``Coming up with features is difficult, time-consuming, requires expert knowledge. Applied machine learning is basically feature engineering.'' --- Andrew Ng
\end{quote}

This book devotes Chapter~\ref{ch:feature-engineering} entirely to this topic.

\section{Exercises}

\begin{exercise}
	For each of the following tasks, state whether it is supervised or unsupervised learning, and whether it is regression or classification:
	\begin{itemize}
		\item Predicting tomorrow's stock price.
		\item Grouping customer purchase histories into segments.
		\item Detecting fraudulent transactions from labeled examples.
		\item Compressing a dataset to 50 dimensions.
	\end{itemize}
\end{exercise}

\begin{exercise}
	A linear model has high bias and low variance. A 15th-degree polynomial has low bias and high variance. Sketch the bias-variance decomposition as model complexity increases. Where is the optimal point?
\end{exercise}

\begin{exercise}
	Suppose the true function is $f^*(x) = \sin(x)$ and we fit a linear model $f(x) = ax + b$. Decompose the expected test error at $x = \pi/2$ into bias and variance components (conceptually; numerical computation not required).
\end{exercise}

\begin{exercise}
	Explain why the No Free Lunch theorem implies that we must use domain knowledge to choose a model class. Provide a concrete example.
\end{exercise}
